\documentclass[12pt]{article}
\usepackage{geometry,fancyhdr,xr,hyperref,ifpdf,amsmath,rcs,indentfirst}
\usepackage{lastpage,longtable,Ventry,url,paunits,shortcuts,smallsec,color,tightlist,float}
\geometry{letterpaper,top=50pt,hmargin={20mm,20mm},headheight=15pt} 
\usepackage[stable]{footmisc}

\pagestyle{fancy} 

\RCS $Revision: 1.5 $
\RCS $Date: 2002/01/09 03:50:54 $

\fancypagestyle{first}{
\lhead{Wet bulb equation}
\chead{}
\rhead{page~\thepage/\pageref{LastPage}}
\lfoot{} 
\cfoot{} 
\rfoot{}
}

\ifpdf
    \usepackage[pdftex]{graphicx} 
    \usepackage{hyperref}
    \pdfcompresslevel=0
    \DeclareGraphicsExtensions{.pdf,.jpg,.mps,.png}
\else
    \usepackage{hyperref}
    \usepackage[dvips]{graphicx}
    \DeclareGraphicsRule{.eps.gz}{eps}{.eps.bb}{`gzip -d #1}
    \DeclareGraphicsExtensions{.eps,.eps.gz}
\fi

\newcommand{\dbar}{d\mkern-6mu\mathchar'26} 
\begin{document}
\newcommand{\vect}[1]{\boldsymbol{\vec{#1}}}
\pagestyle{first}


\section{Introduction}
\label{sec:introduction}

This note expands on Thompkins p. 36 on the wet bulb process and the evaporation of precipitation.

A wet bulb thermometer is a good example of the first law in action.
It is really a pair of thermometers, one measuring
the temperature of incoming  air with vapor mixing ratio $w_v$, and the other the temperature of a saturated
wick evaporating into the air.  As the air blows past the two thermometers,
moist entlapy (as defined in  the
\href{https://drive.google.com/file/d/1BPTvFijtktG4135T1X8DXwI8l_C8xHEt/view?usp=sharing}%
{moist static energy notes} is conserved.


\section{Wet bulb enthalpy}
\label{sec:wet-bulb-enthalpy}


To see how this works, consider the system that consists of 

\begin{itemize}
\item A wet wick with liquid water mass $m_l$ and enthalpy $m_l c_l T_{w}$
\item $m_d$ kg of dry air with enthalpy of $m_d c_{pd} T_{dry}$
\item $m_v$ kg of water vapor with enthalpy $h_v(T_{dry})=c_{pv} T_{dry} + l_{v0}$
\end{itemize}

When the air/vapor hits the wick water will evaporate until the liquid, vapor
and dry air are all at  the new temperature $T_w$ which we know is cooler\
that $T_{dry}$ (since the air coming in is unsaturated).  The new enthalpy
of the sistem will now be the sum of:


\begin{itemize}
\item A wet wick with liquid water mass $m^\prime_l$ and enthalpy $m^\prime_l c_l T_{w}$
\item $m_d$ kg of dry air with enthalpy of $m_d c_{pd} T_{w}$
\item $m_{vs}$ kg of water vapor with enthalpy $h_v(T_{w})=c_{pv} T_{w} + l_{v0}$
\end{itemize}

Since water is conserved we know that the gain in vapor has to be equal to the
loss in liquid:

\begin{equation}
  \label{eq:liquid_bal}
  (m_l - m_l^\prime ) = (m_{vs} - m_v) = m_d ( w_{vs} - w_v )
\end{equation}


So put this all together:

\begin{equation}
  \label{eq:hdry}
H_{before} = m_l c_l T_w + m_d c_{pd} T_{dry} + m_v h_v(T_{dry})
\end{equation}

\begin{equation}
  \label{eq:twet}
H_{after}=  m_l^\prime c_l T_w + m_d c_{pd} T_w + m_{vs} h_v(T_w)
\end{equation}


According to the first law:

\begin{equation}
  \label{eq:first}
  Q dt = dH - V dp
\end{equation}
and since $dp$=0 and $Q$=0, $H$ for the mixture has to stay constant.

Since $H_{before}$=$H_{after}$, I can subtract them and rearrange to get:

\begin{equation}
  \label{eq:subtract}
(m_l - m_l^\prime) c_l T_w + m_d c_{pd} (T_d - T_w)
= m_{vs} h_v(T_w) - m_v h_v(T_{dry})
\end{equation}

Finally, divide through by $m_d$ and  use (\ref{eq:liquid_bal})  and (\ref{eq:subtract}) becomes:

\begin{equation}
  \label{eq:subtract2}
(w_{vs} - w_v) c_l T_w +  c_{pd} (T_d - T_w)
= w_{vs} h_v(T_w) - w_v h_v(T_{dry})
\end{equation}

or remembering that $c_l T = h_l$

\begin{align}
  \label{eq:subtract2}
  c_{pd} (T_d - T_w)
=& -w_v (h_v(T_{dry}) - c_l T_w) + w_{vs} (h_v(T_w) -  c_lT_w) \nonumber\\
  c_{pd} (T_d - T_w) \approx&   -w_v (h_v - h_l) + w_{vs}(T_w) (h_v - h_l)
\end{align}



There's an approximate sign on (\ref{eq:subtract2}) because I've neglected the temperature
difference in the $h_v$  by saying that $c_{pv} T_{dry} + l_{v0} \approx c_{pv} T_{w} + l_{v0}$.
(Question, what is $l_{v0}$ and how big an error am I making here?)  So:

\begin{equation}
  \label{eq:latent}
(h_v(T_{dry}) - c_l T_w ) = (h_v(T_w) -  c_lT_w) = h_v - h_l = l_v
\end{equation}

\section{The wet bulb equation}
\label{sec:wet-bulb-equation}


With that approximation I've got:

\begin{equation}
  \label{eq:wetfinal}
  c_{pd} (T_{dry} - T_w) \approx  l_v ( w_{vs}(T_w) -w_v)
\end{equation}

What is (\ref{eq:wetfinal}) good for?  I can measure $T_{dry}$ and $T_w$, calculate
$w_{vs}(T_w)$ (if I also know the pressure), and solve for the
mixing ratio, $w_v$.

As Thompkins notes -- raindrops also behave like wet bulbs, so that the wet bulb equation is generally important for the evaporative cooling of air by falling precipitation.


\end{document}




%%% Local Variables:
%%% mode: latex
%%% TeX-master: t
%%% End:
